\begin{table}[H]
\centering
\caption{Summary Statistics}
\begin{threeparttable}
\begin{tabular}{lcccc}
\toprule
Variable & Mean & Std.~Dev. & Min & Max \\
\midrule
\multicolumn{5}{l}{\textit{Panel A: Short-Run Panel (Children 8--14 in 1910)}} \\[3pt]
Child labor rate (1910) & 0.090 & 0.098 & 0.014 & 0.355 \\
Child labor rate (1920) & 0.651 & 0.060 & 0.560 & 0.815 \\
School attendance (1910) & 0.943 & 0.057 & 0.767 & 0.987 \\
School attendance (1920) & 0.160 & 0.040 & 0.087 & 0.236 \\
Share male & 0.574 & 0.027 & 0.507 & 0.623 \\
Share white & 0.915 & 0.123 & 0.538 & 0.998 \\
N children per state & 128737.571 & 125243.467 & 3518.000 & 579169.000 \\
[6pt]
\multicolumn{5}{l}{\textit{Panel B: Long-Run Panel (Children 0--10 in 1920, Outcomes in 1940)}} \\[3pt]
SEI (1940) & 22.350 & 3.189 & & \\
Occscore (1940) & 17.159 & 1.705 & & \\
Labor force participation (1940) & 0.790 & 0.020 & & \\
Farm residence (1940) & 0.253 & 0.135 & & \\
Home ownership (1940) & 0.453 & 0.061 & & \\
School attendance (1930) & 0.733 & 0.037 & & \\
MP exposure (years) & 4.184 & 2.855 & & \\
N children per state & 195916.102 & 193648.411 & & \\
\bottomrule
\end{tabular}
\begin{tablenotes}[flushleft]
\small
\item \textit{Notes:} Panel A: State-level aggregates from the IPUMS MLP 1910--1920 linked panel. N = 49 states. Child labor rate = share of children aged 8--14 in the labor force. School attendance = share attending school. Panel B: State-level aggregates from the IPUMS MLP 1920--1930--1940 linked panel. N = 49 states. SEI = Duncan Socioeconomic Index. Occscore = occupational income score. MP exposure = years between mothers' pension adoption and 1920.
\end{tablenotes}
\end{threeparttable}
\label{tab:summary}
\end{table}
